\section{Lecture 24: Cryptography and Security Summary}

\subsection{Overview}
This lecture introduces the core cryptographic primitives and explains how they are combined to build secure communication protocols such as TLS. It focuses on both theoretical foundations and practical applications.

\subsection{Security Requirements}
Four main security goals exist:
\begin{itemize}
    \item \textbf{Confidentiality}: Prevent unauthorized reading of data
    \item \textbf{Integrity}: Ensure data is not altered
    \item \textbf{Authentication}: Verify identity of sender/receiver
    \item \textbf{Non-repudiation}: Prevent denial of sending a message
\end{itemize}

Each requirement maps to a cryptographic primitive. :contentReference[oaicite:0]{index=0}

\subsection{Symmetric Cryptography}
\begin{itemize}
    \item Uses a shared secret key $k$
    \item Encryption: $c = E_k(m)$, Decryption: $m = D_k(c)$
    \item Common algorithm: AES (fast and hardware-accelerated)
\end{itemize}

\textbf{Modes of operation:}
\begin{itemize}
    \item ECB: insecure (leaks patterns)
    \item CBC: provides confidentiality only
    \item GCM: provides confidentiality + integrity (used in TLS 1.3)
\end{itemize}

\subsection{Hash Functions}
A hash maps arbitrary input to fixed output:
\[
h : \{0,1\}^* \rightarrow \{0,1\}^n
\]

\textbf{Key properties:}
\begin{itemize}
    \item Preimage resistance
    \item Second-preimage resistance
    \item Collision resistance
\end{itemize}

Used for integrity checks, password storage, and as building blocks for MACs and signatures. :contentReference[oaicite:1]{index=1}

\subsection{Asymmetric Cryptography (RSA)}
Uses key pairs:
\begin{itemize}
    \item Public key: $(n, e)$
    \item Private key: $(n, d)$
\end{itemize}

\textbf{Encryption:}
\[
c = m^e \mod n
\]

\textbf{Decryption:}
\[
m = c^d \mod n
\]

\textbf{Signing:}
\[
s = h(m)^d \mod n
\]

Provides confidentiality, authentication, integrity, and non-repudiation.

\subsection{Diffie-Hellman Key Exchange}
Allows two parties to establish a shared secret without transmitting it:
\begin{itemize}
    \item Public: $g, p$
    \item Shared secret: $g^{ab} \mod p$
\end{itemize}

\textbf{Forward secrecy}: even if long-term keys are compromised later, past sessions remain secure. :contentReference[oaicite:2]{index=2}

\subsection{Message Authentication Codes (MACs)}
\begin{itemize}
    \item Provide integrity and authentication
    \item Example: HMAC
\end{itemize}

\subsection{Digital Signatures}
\begin{itemize}
    \item Created with private key, verified with public key
    \item Provide integrity, authentication, and non-repudiation
\end{itemize}

\subsection{Public Key Infrastructure (PKI)}
\begin{itemize}
    \item Certificates bind identities to public keys
    \item Verified through a chain of trust (CA $\rightarrow$ intermediate $\rightarrow$ leaf)
\end{itemize}

\subsection{TLS 1.3 Handshake}
Steps:
\begin{enumerate}
    \item Client sends supported ciphers + DH share
    \item Server responds with DH share + certificate + signature
    \item Client verifies certificate and signature
    \item Both derive shared key using DH + hash (HKDF)
    \item Data is encrypted using AES-GCM
\end{enumerate}

\textbf{Key idea:}
\begin{itemize}
    \item Asymmetric crypto: used during handshake
    \item Symmetric crypto: used for data transfer (fast)
\end{itemize}

\subsection{Conclusion}
Modern secure communication relies on combining multiple primitives:
\begin{itemize}
    \item Diffie-Hellman for key exchange
    \item Signatures + PKI for authentication
    \item AES-GCM for encryption and integrity
    \item Hash functions for key derivation and verification
\end{itemize}

TLS integrates all these components to provide secure communication over the internet. :contentReference[oaicite:3]{index=3}